%% kapa: exam=1 round=29 year=2018 period=2 form=A subject=appraisal_law subject_ko=감정평가관계법규 seq=15 number=15
%% source: Q-Net 감정평가사 기출문제 (cst00309) artlSeq=5202808 file=extracted/1008/5202808/2018년 제29회 감정평가사 1차시험 기출문제/2018년 제29회 감정평가사 1차시험 기출문제/2교시 A형.pdf
%% conversion: kordoc 4.24.51 -> tools/gen1.py (auto); review=proofread-sonnet; flags=none
\begin{question}{15}
감정평가 및 감정평가사에 관한 법령상 감정평가업자(감정평가법인 또는 감정평가사사무소의 소속 감정평가사를 포함한다)의 의무에 관한 설명으로 옳은 것은?
\begin{choices}
  \item 자신이 소유한 토지에 대하여 감정평가를 하기 위해서는 국토교통부장관의 허가를 받아야 한다.
  \item 감정평가업무와 관련하여 필요한 경우에만 국토교통부장관의 허가를 받아 토지의 매매업을 직접 할 수 있다.
  \item 감정평가업무의 경쟁력 강화를 위해 필요한 경우 감정평가 수주의 대가로 일정한 재산상의 이익을 제공할 수 있다.
  \item 신의와 성실로써 공정하게 감정평가를 하여야 하며, 고의 또는 중대한 과실로 잘못된 평가를 하여서는 아니된다.
  \item 공익을 위해 필요한 경우에는 다른 사람에게 자기의 자격증을 대여할 수 있다.
\end{choices}
\end{question}
\answer{4}
